\documentclass[exercice,a4]{thedocument}%
\startdatakeys
    \source{laclassecooperativedestef.com}
    \BO{2026}
    \cycle{3}
    \competencesDuSocle{RE4,CO2}
    \objectifsApprentissage{C3C1ab,C3C1aj,C3C1ad,C3C2ac,C3C2ab,C3C1ae}
    \motsCles{perpendicularite, cercle, triangle}
\enddatakeys
\begin{document}%
\noindent\textbf{Programme de construction 5}\par\smallskip\noindent
\begin{itemize}
    \item Tracer le segment $[\mathspoint{AB}]$ tel que
    $\mathspoint{AB}=7\text{ cm}$.
    \item Tracer la droite $(\mathspoint{d})$ perpendiculaire à
    $[\mathspoint{AB}]$ passant par le milieu de $[\mathspoint{AB}]$.
    \item La droite $(\mathspoint{d})$ coupe $[\mathspoint{AB}]$ en
    $\mathspoint{C}$.
    \item Placer le point $\mathspoint{E}$ tel que $\mathspoint{EC}=6\text{ cm}$.
    \item Tracer les segments $[\mathspoint{AE}]$ et $[\mathspoint{EB}]$.
    \item Tracer la droite $(\mathspoint{f})$ passant par $\mathspoint{E}$,
    perpendiculaire à $(\mathspoint{d})$.
    \item Tracer le cercle $\mathcal{C}$ de centre $\mathspoint{E}$ et de
    rayon $3{,}5\text{ cm}$.
    \item Le cercle $\mathcal{C}$ coupe la droite $(\mathspoint{f})$ en
    $\mathspoint{G}$ et $\mathspoint{H}$.
    \item Tracer les segments $[\mathspoint{AG}]$ et $[\mathspoint{BH}]$.
\end{itemize}
\end{document}%
